\section{Karbala}

\begin{quotex}
These extracts from the longer work \emph{Persian Traditions in Spain}\footnote{\url{https://www.gornahoor.net/?page\_id=11751}} by \textbf{Michael McClain} were originally published in the
\emph{International Journal of Shi'i Studies}, Volume V, No. 1, 2007, under the title ``Affinities Between Catholicism and Shi'ism''.
\end{quotex}

The special affinity\footnote{Excerpted from \emph{Persian Traditions in Spain} by Michael McClain.} between traditional Catholicism and Shi`ism has been noted by many, indeed since the time of the Crusades. In the twelfth century, the Crusader Archbishop William of Tyre wrote: ``Muhammad's son-in-law `Ali (Ibn Abi Talib) was the best knight, braver and more valiant than any of the other caliphs had been... The Shi'a is not so far from the true Christian faith as is the Sunna.''\footnote{Zoe Oldenbourg, \emph{The Crusades} (New York, 1966), p. 492.} Far more recently, the same sentiments were echoed in \emph{Roman Catholics and Shi`i Muslims}, by James A. Bill and John Alden Williams (Chapel Hill, North Carolina, 2002).

As we shall see below, the great Spanish Catholic poet and mystic \textbf{St. John of the Cross} could in a very real sense be considered a spiritual master in the line of Hassan, the second Shi`a Imam and the son of 'Ali Ibn Abi Talib.

The above is confirmed by the fact that these feelings are very much mutual. Says \textbf{Mahmoud Ayoub}, a Lebanese Shi'a:

\begin{quotex}
The problem of human freedom and divine sovereignty and will is as complex a theological issue in Islam as it is in Christianity. While both the broad emphasis and the more basic theological orientation are quite different in the two traditions, the Shi'i view is closer to the Christian position than it is to the strict Sunni Islamic view. God wills, knows and decrees; yet man is still responsible for his choice, a choice which confronts him at every moment, as the earth would never be void of a proof or witness (hujjah) of God over His servants both to judge and to redeem them. The proximity of the Shi`i view to that of Christianity is perhaps due to the fact that both accept a mediator between man and God, one whose essential being and place in human history plays a determining role in the divine plan for creation, revelation and salvation. Thus we must agree with Henry Corbin that Shi`i imamology is a kind of ``Islamic Christology''. In Christian piety, (Jesus) Christ is the eternal Logos, the divine Word; the agent of creation on one hand, and on the other hand the slain lamb standing before the throne of majesty both to save and to judge. The Imams, likewise, are at one and the same time the pivot of creation and reason for its subsistence, and the blood-stained martyrs whose death is a point of contention between God and their persecutors.
\flright{\textsc{Mahmoud Ayoub}, \emph{Redemptive Suffering in Islam: A Study of the Devotional Aspects of Ashura in Twelver Shi'ism}, The Hague, 1978, p. 199.}
\end{quotex}


\textbf{Sayyid Musa al-Sadr}, the Iranian-born former spiritual leader of the Lebanese Shi'as, illustrates this. He sometimes preached in Catholic churches, and showed himself a past master in touching Catholic hearts and evoking Catholic spirituality, very often moving the congregation to tears, as in his famous Lenten homily of February 19, 1975 in the Cathedrale Saint-Louis des Capucins in Beirut.\footnote{See \textsc{Fouad Ajami}, \emph{The Vanished Imam} (Ithaca, New York, 1987), pp. 133-135. \textsc{Majed Halawi}, \emph{A Lebanon Defied: Musa al-Sadr and the Shi'a Community} (Oxford, England, 1992), pp. 192-195.} Though Sayyid Musa al-Sadr was certainly a brilliant man, this above-mentioned mastery, much greater than that of most Catholic priests who have studied for years at a Catholic seminary and have long experience in the Catholic priesthood, would be inconceivable if his training at Shi'ite centers in Qum and Najaf had not somehow prepared him to touch Catholic hearts and evoke Catholic spirituality.

Traditional Catholics are much closer to traditional Shi'as than they are to either Protestants or ``modernist,'' ``progressive'' and ``post-Vatican II'' Catholics. The Catholic thinker Hilaire Belloc considered Protestantism to be the fountainhead of modernity, and therefore something very near to ``the source of all evils.'' However, he had considerable respect for Islam, and left no doubt that he believed Islam to be far closer to Catholicism than is Protestantism.\footnote{Hilaire Belloc, \emph{The Crisis of Civilization}, New York, 1937.}

Both traditional (Catholic and Eastern Orthodox) Christianity and traditional Islam affirm that God is both immanent and transcendent. To deny either God's transcendence or His immanence is to put limitations and conditions on God. Besides the general affinity between traditional Catholicism and Shi`ism, Spanish Catholicism and Iranian Shi`ism have so many particular characteristics in common that the resemblance between them is often uncanny. We will give only a few examples. The resemblance between Ashura, the anniversary of the martyrdom of Imam Hussayn as celebrated in Iran, and Holy Week as celebrated in most of Spain has been noted by many. Of course, the Church very strongly disapproves of deliberately doing bodily harm to oneself. However, anyone who has heard the heart-wrenching ``saetas'' of a Spanish Holy Week and seen the processions with the images of Christ being scourged at the pillar and carrying the Cross will readily grasp what we are talking about. During Holy Week in Spain, I have heard keening and seen weeping. I have seen people prostrate themselves, carry or drag heavy weights long distances or hobble long distances on their knees. In some places in Spain, notably in Lorca in the province of Murcia, dramatic representations of the events leading to the Crucifixion are part of the Holy Week ceremonies, demonstrating a striking parallel to the \emph{taziyeh} of Ashura celebrated in Iran...

Below is a poem by the Persian poet \textbf{Ansari} on the martyrdom of \textbf{Imam Hussayn}:

\begin{verse}
O breeze of morning, take to 'Ali these words of the poet Ansari:\\
Say: Hussayn is fallen. Rise, then, go and see.\\
To Karbala from Najaf where you lie,\\
His body in a hundred pieces pierced by the lance, the dagger, the sword.\\
See who was once the light of your eyes,\\
Now the enemy around him like eyelashes around the eye;\\
And here you lie, in pleasant repose with Adam and Noah, at rest\\
While Hussayn has as his resting place the burning sands of Karbala!\\
Although you were made stranger to yourself by the stroke of the sword,\\
Around you were both stranger and kin, with refreshments and sweets;\\
While the body of your (son) Hussayn is rent the whole length with wounds.\\
And would you know the number of those wounds?\\
They are as many as the stars!\\
Wherever you turned your gaze, there stood a friend to see,\\
While Hussayn's eye falls only on the enemy.\\
`Ali, when you gave your life your family was there beside you,\\
But there on a desert plain far from daughter or sister Hussayn dies.\\
Faithful Spirit, Gabriel, brought a shroud from Heaven,\\
But Hussayn fell there on the earth without ablution, without shroud!\\
`Ali, since Hussayn in the last hour took your head on his lap to lie,\\
As kindness in return, then, lay his head on your lap till he dies.\\
\flright{\textsc{Lynda Clark}, \emph{Elegy (Marthiya) on Hussein: Arabic and Persian}, Al-Serat, Spring \& Autumn (London 1986), pp. 24-25.}
\end{verse}


According to Ibn Abbas, after the battle of Siffin, '\textbf{Ali Ibn Abi Talib}, the First Imam, told him:

\begin{quotex}
Jesus one day passed with his disciples through Karbala and on that spot they saw a group of gazelles gathered together weeping. Jesus and his disciples sat and wept with them, without the disciples knowing the reason for that lamentation. Jesus finally told them that this was a spot on which was to be killed the young descendant of the Apostle Ahmad (Muhammad), and child of the pure, unblemished virgin Fatima, who is like unto my mother (the Virgin Mary). He (Hussayn) shall be buried on this spot whose soil is more fragrant than musk. For it is the burial place of the martyr (Hussayn). Such is the soil containing the bodies of prophets and descendants of prophets.
\flright{\textsc{Ayoub}, op. cit., pp. 237-238.}
\end{quotex}


Some will say that `Ali Ibn Abi Talib, the First Imam, was trying to express a difficult truth by way of allegory and parable, as Jesus often did. However, at this point, there is something I wish to note. In Arabic, the New Testament or the Gospel is called \emph{Injil}, which is not pure Arabic, but rather it is derived from the Greek \emph{Euangelion}, meaning, roughly, ``good news'' or ``a good message''. Like the word ``angel,'' the Greek \emph{Euangelion} is remotely derived from the Avestan word for ``messenger,'' the \emph{Eu} being a prefix which means ``good''. From the Greek \emph{Euangelion} comes the Medieval Latin \emph{Evangelium}, from which in turn are derived so many words in English and various Romance languages...

We shall deal later with traditions connecting Imam Hussayn and \textbf{St. John the Baptist}. For the moment, we continue to refer to the tragedy of Karbala and the martyrdom of Imam Hussayn. Below is part of a meditation on the month of Muharram, which marks the anniversary of the tragedy of Karbala. This one is particularly interesting, because it includes a brief account of the tragedy written by \textbf{Ja`far al-Sadiq}, the Sixth Imam and great-grandson of the Imam Hussayn, the martyr of Karbala.

\begin{quotex}
Muharram is one of the four months declared sacred by Allah in the Holy Qur'an. The other three holy months are Rajab, Dhu Qi`da, and Dhu Hijja. ``There is no god save Allah.'' If someone claims to be a god or has faith in others' godhead, other than Allah, Allah does not make a compromise to take anyone as god, although He gives a free hand even to such false pretenders, because: ``There is no compulsion in religion'' (Qur'an: 2:156).
\end{quotex}


In verse 59 of Sura al-Nissa, Allah says: ``O you who believe! Obey Allah, and obey the Messenger and the \emph{Ulil Amr} (those who are authorized to command) from among you.'' The believers always obey Allah, His messenger, and the \emph{Ulil Amr} appointed by Him. It is an open invitation. There is no compulsion. Those who are not believers can do what they wish. Likewise the prophets and messengers sent down by Allah to guide and warn people never accepted anyone else as a prophet or messenger, whatever the circumstances, although they did not stop claimants of this kind, or opponents, with the application of force or repression.

In the same manner, Imam `Ali Ibn Abi Talib did not recognize anyone else as \emph{Ulil Amr}. His elder son Hassan, like his grandfather, the Holy Prophet (with the treaty of Hudaybiah), let Mua`wiyah become the ruler, but he did not give up his rights as \emph{Ulil Amr}, which he was at that time.

When Yazid became the ruler, he threw to the wind the earlier policy of the rulers not to demand the \emph{bay`a} (oath of obedience) from the children of the Holy Prophet, and he began to exercise pressure upon the Imam Hussayn to swear fealty and acknowledge him as the \emph{Ulil Amr}, which the Holy Imam rightly refused to do. Another aspect of the whole affair is that Yazid was the vilest tyrant of the worst type, but even if he had been an ordinary ruler Imam Hussayn could not swear fealty to him as the \emph{Ulil Amr}. So he did not. In the month of Rajab, in year 68 A.H., he left Medina and went to Mecca. From Mecca, before performing the Hajj, he took his family with him, and with some friends and companions, he began the journey toward Iraq, with the expressed intention to cross the boundaries of the empire under the rule of Yazid, and settle down in some other country, Iran or India, because he wanted to make it clear to Yazid that he could not swear obedience to a non-\emph{Ulil Amr}, as he himself was an \emph{Ulil Amr}.

It was not to be. A large army in service of Yazid under the command of `Umar Ibn Saad surrounded the caravan of the Imam Hussayn when he reached Karbala on the second of Muharram, 61 A.H.

There are several aspects of human relationship and behavior in the events of Karbala that took place in the 10 days of Muharram, culminating in the Martyrdom of the Imam Hussayn and his 72 friends and relatives, which come into sharp focus as we recount every minute, during the religious gatherings each year, but above all, the ultimate reason remains the same: the impossibility of swearing an ``oath of loyalty'' (\emph{bay`a}) by an \emph{Ulil Amr} to obey a non-\emph{Ulil Amr}.

After the Imam Hussayn, all his successors to ``the office of \emph{Ulil Amr},'' our Holy Imams, refused to obey ``non-\emph{Ulil Amrs},'' and every ruler held each of them prisoner, used every trick, applied force, and in the end killed every Imam, exactly as Yazid did.

Muharram is a month of mourning for the lovers and followers of the family of Muhammad. On the tenth day of this month in 61 A.H., the Imam Hussayn, the grandson of the Holy Prophet and the younger son of 'Ali and Fatima, together with his family and friends, in all 72 men, were slain on the sands of the desert of Karbala. Since then, each year the true followers of the Holy Prophet, through grief, sorrow and tears, keep alive the message, cause and purpose of the greatest martyrdom in human history.

No doubt Muharram is a holy and sacred month. The believing men and women, in this month, suspend the application of good effects of days and dates and avoid rejoicing even if happy events come to pass. The friends and followers of the family of Muhammad hold meetings; they have been doing so for the last 1,300 years, in the name of the Imam Hussayn, during the months of Muharram and Safar, particularly in the last ten days of Muharram, to give new life to the Divine Message of ``There is no god but God.''... Each year at the advent of Muharram, Islam turns over a new leaf. In fact, it is on account of the Imam Hussayn's remembrance every year that we know who the Holy Prophet, `Ali Ibn Abi Talib, and Bibi Fatima Zahra are, and what their true substance, style and wisdom were. It is because of the Imam Hussayn's memory that we call to mind each year the original Islam, and it becomes clearly visible from behind the smoke screen of the dust of delusion thrown into the eyes of Muslims in the name of Muslim rule... .

Ja`far al-Saddiq, the Sixth Imam, recounts:

\begin{quotationx}
A large army in the service of Yazid, under the command of `Umar Ibn Saad, surrounded the Imam Hussayn, his family and friends, on all sides on the ninth of Muharram, shouting for joy and uttering songs of triumph, in view of the presence of countless soldiers armed with swords, lances, arrows, and sling stones, all this against a handful of 72, among whom were 90-year-old men and a six-month old child. The Imam Hussayn at that time received no support from any quarter. He was forsaken in this dire situation.

On the night of Ashura, the Holy Imam had invited his companions and relatives to sit together, after evening prayers, and hear his last address. As soon as the prayers came to an end, they sat around him.

He said: ``Brothers! It is a misjudgment if, in one's heart of hearts, anyone expects to witness my victory in tomorrow's battle and gather the fruit of conquest. I tell you, in clear words, that the enemy will slay me, beyond the shadow of a doubt, and in a manner that will make your blood run cold. `Abbas, my brother, his hands both cut off, shall be killed on the slope of the shore of the River Euphrates, running along the river. A lance will pass through the heart of my son, `Ali Akbar. The hooves of the enemy's cavalry will crush the wounded body of my brother's son, Qasim. Likewise, whoever stays with me, be he a kinsman or a friend, will be killed in cold blood. I tell you, beforehand, that an arrow, deliberately and purposefully aimed with premeditated malice, shall hit the throat of my six-month old son, `Ali Asghar. He shall die on the spot in my arms. My son, `Ali, who now suffers from a high fever, alone shall weather the storm, and bear in the aftermath the worst and meanest abuse, along with the womenfolk of Muhammad's family.

``Friends! Do not let your heart break nor have qualms on account of obligation of loyalty to me and my cause, concerning which you have sworn an oath. I, willingly, set you free from the sworn pledge. Leave me alone to that which lies in wait for me tomorrow. They will be after me and no other. The night is dark. You can escape and vanish unseen in the darkness.''

The the Holy Imam asked his brother `Abbas to put out the candles, to let the deserters show their heels and save their necks. In the darkness many escaped. The lamps again were lit and a final group of 72 was sitting there calm and quiet, willingly prepared to meet certain death the next morning when the ``great sacrifice'' (\emph{dhibh `azeem}), as promised by the Almighty, would become a fact.

``And We (Allah) ransomed him (Ismail) with a great sacrifice'' (Qur'an, 37:107). It was a dark, dismal and deadly night. The chosen group of godly men and women passed the whole night reciting the Holy Qur'an, worshipping Allah, performing Sal?t, and their guide, the Imam Hussayn, was with them at all times.

It is a proof of sincere devotion to Allah if we, his followers, also spend this night in the worship of Allah. Whoever spends this night awake, until the daybreak, worshipping Allah, near the grave of the Imam Hussayn, Allah will raise him or her on the Day of Judgment, along with the martyrs of Karbala...
\end{quotationx}

Today is the Day of Ashura, the day of Martyrdom of Imam Hussayn. Today, in the year 61 A.H., the Imam the S?ra al-S?ff?t: ``And We (Allah) ransomed him (Ismail) with a great sacrifice.''

It was the third day since the water supply to the Imam's camp had been disconnected. At dawn the small devoted group performed ``Tayammum'' and assembled to perform the Fajr Sal?t behind the Imam Hussayn. The first swarming flight of arrows shot by the enemy archers hit the Holy Imam and his devotees while they were saying the Sal?m. Thirty persons were killed on the spot. With ``whatever is left,'' the Imam came onto the battlefield and before attacking the enemy soldiers, addressed them to make known his rights and status with reference to the Holy Prophet and the Book of Allah. He warned them not to kill him as it would bring upon them destruction in this world and everlasting punishment on the Day of Judgment. The Holy Imam again and again tried to make them understand the consequences of the events that would take place if they did not listen to him. He once again proposed that they let him, together with his family and friends, go away and settle down in some far-off land, far from the jurisdiction of the Muslim Empire. They listened to what he had to say, but declared: ``We shall kill you if you do not agree to recognize Yazid's overlordship by swearing loyalty to him.''

Despite the hunger, thirst, and wounds, one by one the devotees of Imam Hussayn went to fight against the hordes of demons in human shape, displayed rare acts of bravery and courage, and gave their lives in the cause of Allah. On every occasion the Holy Imam, together with his brother and son, dispersed the ``blood hounds,'' carried away the dead bodies, and lay down on the ground under a tent, now known as ``Ganj-i-Shaheedan''.

The sun had crossed the meridian. The time for the Zuhr prayers had begun. The renegades refused any sort of truce. The Imam therefore prayed Zuhr Sal?t as ``Namaz of Khawf''.

After the companions and friends of the Holy Prophet, the Imam Ali, and the Imam Hussayn had given their lives, as he had foretold in his ``Shab-i-Ashura'' speech, it was the turn of his relatives. The two sons of Bibi Zaynab (Awn, 10, and Muhammad, 9), Qasim, the son of Imam Hassan, 14, `Abbas, the standard bearer, the backbone of the Imam, `Ali Akbar, the Imam's 18 year old son, the other sons of `Ali Ibn Abi Talib, and the grandsons of `Ali Ibn Abi Talib, one and all, gave their lives in order to bring to life the message ``there is no god but God'' and for keeping alive their Imam, \emph{Ulil Amr}.

Then our beloved Imam stood alone. Bruised, slashed, cut, gashed, thirsty, soaked in his own blood and in the blood of his sons, brothers, nephews and devoted friends, he went into the tent of his sister, Bibi Zaynab. His six-month old son, `Ali Asghar, was in her lap, dying of thirst and hunger. He took him in his arms and slowly walked up to the warlike array of heartless Muslims.

He lifted the child up on his hands. This sad, sensitive, eloquent spectacle made even the devil's disciples curse the devil. With tears in their eyes they cried in deep anguish. Afraid of a revolt in the ranks, Ibn Saad looked at Hurmalah, who, reading the message in his commander's eyes, took aim and shot the fatal arrow which, passing through the Imam's hand, went into the child's neck. `Ali Asghar stared at his father's face, smiled, and rested in peace. The Imam dug a small grave and buried his son.

Then the Holy Imam paid his last visit to his kinfolk in the tents and went forth mounted on his horse Zuljenah, wearing his grandfather's garments, and holding ``Zulfiqar'' (the sword of `Ali Ibn Abi Talib), in his hand.

On the battlefield he pronounced his last call: ``Is there a helper?'' There was no response. A fierce battle ensued. Then Zuljenah sat on the ground. The Holy Imam, every pore of his body a bleeding wound, slid over the burning sands of Karbala. It was the time for the `As?r prayers. He bent over and rested his forehead on the ground in prostration. Swords, arrows, spears, daggers, lances, and stones hit him from all directions. The sacred blood of Muhammad, `Ali Ibn Abi Talib, and Fatima Zahra flowed in a stream over the sandy soil.

The Imam Hussayn, our master, and master of every living being, thereupon whispered to his Creator: ``O merciful Lord of the worlds! Hussayn, your servant, has given up everything which you had granted him according to your holy will. Accept the humble sacrifice of your servant Hussayn. If I, the grandson of your Messenger, had more, I would have surrendered it to you willingly... .''

The eternally cursed Shimr came close and severed the sacred head of the Holy Imam from his blessed body. The heavens and the earth sank low into a pit of gloom. A dreary, dismal darkness spread out everywhere. A cosmic cry of agony echoed in every nook and cranny of the universe. Animals stopped in their tracks, birds swerved in their flight, and a shiver ran through trees, water, valleys, plains and mountains.

The demons, in devilish frenzy, trampled the bodies of the devoted martyrs under the hooves of their horses, plundered, looted, and set fire to the tents. The bewildered children and tearful women ran to Bibi Zaynab, the daughter of ``Asadullah,'' and gathered around her.

From the first of Muharram to the ninth we discuss the philosophy of Hussayn's martyrdom, we look into and make clear every aspect of true Islam, we carefully identify the role and merits of the Holy Prophet, the Imam `Ali Ibn Abi Talib, Bibi Fatima Zahra, we make known and expound upon the real meanings of the ``Word of Allah,'' the Holy Qur'an, on the strength of the wisdom and knowledge of Muhammad and the Family of Muhammad, we recount all the facets of the character of the Imam Hussayn which achieves much toward making available to mankind the vision and the ideal of making efforts to establish a society free from submission to idols and false gods -- in whatever form, idea or institution they try to enslave mankind -- and obedient to Allah and His laws.

But today we the friends and followers of the Imam Hussayn, who love him and belong to him, and through him to his father, mother and grandfather, and through them belong to their Lord, the Lord of the worlds (1) pray the Fajr Sal?t. After the Fajr Sal?t we carry out the A`m?l Ashura. (2) We do not drink water, nor eat food. (3) Barefoot, bareheaded, attired in black, we come out from the shade and shelter of our homes, in the open, under the sky. (4) Together, disciplined, row after row, we recite moment to moment the events that took place in Karbala, until evening. (5) We pray the Zuhr and `As?r Sal?t at their fixed times. (6) Tears come to our eyes as the dry yellow leaves fall to the ground in autumn. It is an instinctive reaction; we weep, mourn and cry; in the height of our love for the Holy Imam some of us scrape and cut ourselves (not others) to feel the pain and hurt that the martyrs of Karbala suffered on the tenth of Muharram, 61 A.H. In no way is there any hint of self-torture or exhibitionism, just as Uways Qaranee broke all his teeth when he heard the news that in the battle of Uhad his ``beloved'', the Holy Prophet had lost a couple of teeth. It is the domain of love and passionate attachment, deep and abounding devotion, like the leap which the prophet Ibrahim took into the blazing fire. (7) We invoke Allah to lay a curse upon the killers of the Imam Hussayn, his family and friends.''\footnote{Abbas Qummi, \emph{Supplications, Prayers \& Ziyarats} (Iran 2005), pp. 182-195.}


\flrightit{Posted on 2016-08-30 by Michael McClain}

\begin{center}* * *\end{center}

\begin{footnotesize}
\begin{sffamily}
\texttt{Akira Kalashnikov on 2016-09-02 at 19:23 said:}

Fascinating post.

(Okay this is probably going to go straight to the spam bin, but here goes... )

Some further reading:

On the Alawites, the most mysterious of the Shiites, and the closest of all to Christianity:

The Asian Mystery

\url{https://books.google.com/books?id=FBZdAAAAcAAJ\&dq=inauthor\%3A\%22Samuel\%20Lyde\%22&pg=PR3\#v=onepage\&q\&f=false}


The Ansyreeh and Ismaeleeh: A Visit to the Secret Sects of Northern Syria

\url{https://books.google.com/books?id=olMBAAAAQAAJ\&dq=secret\%20sects\%20of\%20syria\&pg=PR1\#v=onepage\&q\&f=false}


SECRET SECTS OF SYRIA

\url{https://archive.org/stream/secretsectsofsyr032392mbp/secretsectsofsyr032392mbp\_djvu.txt}


On the sharing of Shrines between Christians and Muslims:

\url{https://archive.org/details/Hasluck1929ChristianityIslam01}

\url{https://archive.org/details/Hasluck1929ChristianityIslam02}


`On Some Parallels Between Anatolian and Balkan Islamic Heterodox Traditions and the Problem of their Co-Existence and Interchange with Popular Christianity'

\url{https://books.google.com/books?id=i9E29RZgHDIC\&lpg=PA75\&ots=FCr4bVwh4C\&dq=parallels\%20between\%20anatolian\%20balkan\%20heterodox\%20islamic\%20and\%20christian\&pg=PA75\#v=onepage\&q\&f=false}


\hfill

\texttt{Cologero on 2016-10-09 at 08:29 said:}

It seems that Mr. McClain may have been years ahead of his time as this piece \textit{Shiite Catholic}\footnote{\url{https://sancrucensis.wordpress.com/2016/09/30/shiite-catholic/}}, by a Cistercian monk, shows.

The \textit{Vienna-Qom Circle for Catholic-Shi'a Dialogue}\footnote{\url{https://viqocircle.org/}} has been established (it includes Russian Orthodox members) to discuss Religion, Philosophy and Political Theory.

\hfill

\texttt{Daniel on 2016-11-04 at 21:14 said:}

The Persian poet Ansari (Ansari of Herat) was not a Shi'i. He was a Sunni jurist of the Hanbali madhhab and a Sufi disciple of the famous master Abu'l Hasan Kharaqani.

There is some degree of affinity between Shi'ism and Catholicism, but there is a greater affinity between Catholicism and Sufism, as a study of Louis Massignon would show.


\hfill

\end{sffamily}
\end{footnotesize}

